\subsection{Tropical motivic Fubini theorem}\label{jp:fubini}
In this subsection we work in the field $\widetilde{\C}$ of complex
Puiseux series, considered as an elementary substructure of
$\mathbb U$. Put $F=\C((t))$ and take the parameter base to be
$\mathbb S=F$. Normalize the valuation by $v(t)=1$, so that the
value group of $\widetilde{\C}$ is identified with $\Q$.
All definability and all categories are interpreted in the original
language $\lan{}{RV}{}$ over $\mathbb S$. In particular, each
$a/N\in\Q$ is definable over $\mathbb S$ as the unique solution of
$N\gamma=a\,v(t)$.

\begin{rem}[The finite Kummer actions]\label{puiseux:kummer:action}
There are natural identifications
\[
 \gal(\widetilde{\C}/F)\ \cong\ \widehat{\mu}
 \ \cong\ \aut\bigl(\RV(\widetilde{\C})/\RV(F)\bigr),
 \qquad
 \widehat{\mu}=\varprojlim_N\mu_N,
\]
where the transition map $\mu_M\to\mu_N$ for $N\mid M$
is $\zeta\mapsto\zeta^{M/N}$, and the last automorphism group
uses the induced original-language structure on $\RV$.

For the proof, fix $N\geq1$, choose any root $\tau_N$ of
$X^N-t$ in $\widetilde{\C}$, and put
$F_N=\C((\tau_N))$. Grouping exponents modulo $N$ writes each
element of $F_N$ uniquely as
\[
 \sum_{j=0}^{N-1}\tau_N^j f_j(t),\qquad f_j(t)\in F.
\]
Uniqueness also follows from the pairwise distinct valuations
modulo $\Z$ of the nonzero summands. Thus $[F_N:F]=N$.
Its $F$-automorphisms are precisely
$\tau_N\mapsto\zeta\tau_N$ for $\zeta\in\mu_N$.
Replacing $\tau_N$ by another root multiplies it by an element
of $\mu_N\subseteq F$, so both $F_N$ and the ratio
$\sigma(\tau_N)/\tau_N$ are independent of that choice.
For $N\mid M$ these ratios satisfy
\[
 \frac{\sigma(\tau_N)}{\tau_N}
 =
 \left(\frac{\sigma(\tau_M)}{\tau_M}\right)^{M/N},
\]
since $\tau_M^{M/N}/\tau_N\in\mu_N\subseteq F$.
The fields $F_N$ have union $\widetilde{\C}$.
Their finite automorphisms therefore give exactly the compatible
systems $(\zeta_N)_N\in\widehat{\mu}$, proving the first
identification. Each such automorphism preserves the valuation:
in the displayed expansion it multiplies the $j$th summand by
$\zeta_N^j$, preserving its valuation and the unique least
valuation of a nonzero sum.

Write $r_N=\rv(\tau_N)$ and identify $\C^\times$ with the
value-zero subgroup of $\RV(\widetilde{\C})$.
For any $u$ of value $a/N$ there is a unique $c\in\C^\times$
such that $u=c\,r_N^a$. The induced action is
\[
 \sigma_{(\zeta_N)}(c\,r_N^a)
   =c\,\zeta_N^a r_N^a .
\]
In particular, a nonidentity system moves $r_N$ for some $N$.
The action on $\RV$ is faithful.

Finally, $\RV(F)$ is the subgroup generated by
$\C^\times$ and $\rv(t)$, by taking the leading term of a
Laurent series. An automorphism of the induced RV structure
fixing $\RV(F)$ must send $r_N$ to
$\zeta_Nr_N$ for some $\zeta_N\in\mu_N$: these are precisely
the $N$ roots of $\rv(t)$ in $\RV$.
The same relations as above force compatibility of the
$\zeta_N$. The displayed action then determines the
automorphism on every RV element. It agrees with the unique
valued-field automorphism constructed from that system.
This proves the second identification.
The parameter base for all these definability assertions is
still $\mathbb S=F$.
\end{rem}

% BEGIN REFEREE INSERTION R081
\ResolvedRefereeRemark{081}{Clarification; revised}{Normalization and language in the Puiseux setting}{%
\textit{Original source line 1796.}\quad For the Puiseux computation, explicitly normalize $\rcsn(1)=\rv(t)$ and choose compatible roots. Retain the distinction already introduced above: unadorned categories use the original language over $\mathbb C((t))$, while daggered categories use the expansion. A normalized reduced section makes the compatible roots of $t$ definable in the expanded language, so this distinction is necessary to retain monodromy.
}{The author approved the replacement on November 13, 2026. The Puiseux convention now uses only the original RV language over $\mathbb S=\C((t))$. The finite Kummer proof retains monodromy and leaves the auxiliary roots outside the definability base.}
% END REFEREE INSERTION R081

% BEGIN REFEREE INSERTION R082
\ResolvedRefereeRemark{082}{Typo}{Formula for a reduced section}{%
\textit{Original source line 1799.}\quad The assignment should be $1/n\mapsto\rv(\xi_n t^{1/n})$; the present $n\mapsto\xi_nt^{1/n}$ has the wrong input and takes values in $\VF$, rather than $\RV$. Equivalently, write $\rcsn_\xi(a/n)=\rv(\xi_n^a t^{a/n})$ and mention the compatibility that makes this independent of the fraction.
}{The author approved deletion of the reduced-section parametrization. Remark~\ref{puiseux:kummer:action} gives the correctly typed finite-stage RV action, its root-choice independence, and its compatibility under divisibility.}
% END REFEREE INSERTION R082

